\chapter{Metodologia: Estudo de Caso}
\label{ape:metodologia_estudo_caso}

Um estudo de caso é uma metodologia de pesquisa ou análise aprofundada que examina uma situação, problema, indivíduo, grupo ou organização real em seu contexto específico, com o objetivo de compreender fenômenos complexos, identificar causas, consequências e propor soluções ou aprendizados \cite{yin2015}.

Ele é amplamente utilizado em diversas áreas, como administração, educação, engenharia e ciência da computação, porque permite investigar detalhes e particularidades que métodos quantitativos gerais não capturam. Diferentemente de pesquisas apenas estatísticas, o estudo de caso combina coleta de dados qualitativos e quantitativos para gerar uma compreensão contextualizada e profunda do objeto analisado. Em essência, é uma ferramenta que traduz teoria em prática, conectando conceitos acadêmicos a situações reais.